% Yeban User Guide — Chapter 6: Automation and Modulation (English)
% T-DOC-001 chapter-by-chapter delivery, written ahead of the CH06 task.
% Sources: yeban-model/src/project.rs AutomationWriteMode, bridge.rs automation lanes,
% cli.rs SHORTCUTS, yeban-dsp/src/oscillator.rs global LFO, docs/manual/05-mixer-effects.md.

\chapter{Automation and Modulation}
\label{ch:automation-modulation}

\classquote{易穷则变，变则通，通则久。}{Book of Changes, “Appended Judgments II” (《周易·系辞下》)}{Translation: “When a path is exhausted, it changes; through change, it flows; through flow, it endures.” Automation lets parameters change over time; with change, the sound stays alive.}





Automation lets parameters change over time: draw a curve, or record your movements
while playing back. This chapter covers showing lanes, what the four write modes mean,
and the real boundary of modulation in the current build.

\section{Automation Curves and Write Modes}

\subsection{Show the lanes}

Press A in the arrangement area, or use the View automation switch
(Cmd/Ctrl+Shift+A) to show automation lanes. Each non-master track can have its own
lanes, arranged by target parameter. A lane holds a polyline whose vertices are
connected in time order; its horizontal axis aligns with the arrangement timeline, so
you can read the curve directly against clip positions.

\screenshot{ch06-automation-lanes}{Automation lanes expanded in the arrangement area, showing a parameter curve with vertices}

\subsection{Draw and edit the curve}

A lane's polyline is made of vertices. Adding one gives the parameter a value at a
point in time; neighbouring vertices are interpolated. Drag a vertex to change its
time or value; delete one and the curve re-connects across the gap. Curves are saved
with the project, drive parameters during playback and take effect on export too.

\subsection{The four write modes}

Yeban's model defines four automation write modes; knowing them avoids recording
curves you did not mean to:

\begin{description}
  \item[Off] The default. Nothing is written during playback; manual parameter changes do not alter the curve.
  \item[Write] Once the parameter is changed during playback, writing continues until you explicitly stop.
  \item[Touch] Writes only while you are physically touching the parameter; releasing stops immediately. Good for a single corrective nudge.
  \item[Latch] Once touched, writing continues even after you release, until explicitly stopped. Good for locking in a whole fader move.
\end{description}

\begin{table}[htbp]
\centering
\caption{Write modes compared}
\label{tab:ch06-write-modes}
\begin{tabular}{ll}
\hline
Mode & After you release \\
\hline
Off & No writing \\
Write & Keeps writing (explicit stop) \\
Touch & Stops immediately \\
Latch & Keeps writing (explicit stop) \\
\hline
\end{tabular}
\end{table}

\section{Parameter Binding and LFO}

\subsection{Bind a parameter}

An automation lane binds to a specific parameter on a specific track: choose the
target, then draw on that lane. Device parameters (e.g.\ on an insert effect) can also
be automation targets; insert the device, then pick its parameter in the lane's target
list. Parameter changes and automation writes share one commit path, so undo restores
them step by step — you will not get a curve change that the project failed to record.

\subsection{The real LFO boundary}

Yeban's DSP core does implement a \textbf{global LFO} (low-frequency oscillator): one
shared by all voices, evaluated per block, with selectable waveforms. But note
carefully: it is \textbf{not yet wired into the synthesizer's modulation inputs}.
Do not expect to turn an LFO knob in the GUI and get vibrato or filter sweep; those
modulation paths are still evolving.

\begin{figure}[htbp]
\centering
\begin{tikzpicture}[node distance=1.1cm, >=stealth,
    box/.style={draw, rounded corners, align=center, minimum width=2.3cm, minimum height=0.8cm}]
  \node[box] (a) {Pick target};
  \node[box, right=of a] (b) {Pick write mode};
  \node[box, right=of b] (c) {Draw / record};
  \node[box, right=of c] (d) {Playback drives param};
  \node[box, below=of c] (e) {LFO: in DSP, not wired to synth};
  \draw[->] (a)--(b); \draw[->] (b)--(c); \draw[->] (c)--(d);
  \draw[dashed,->] (c)--(e);
\end{tikzpicture}
\caption{Automation flow and the LFO boundary: solid = available path, dashed = modulation not yet wired}
\label{fig:ch06-automation-flow}
\end{figure}

For that reason this chapter does not describe vibrato or filter sweep as available
features. If you cannot find the control in the GUI, it is not wired up yet — you have
not missed it.

\section{Help Topic Mapping}

\begin{table}[htbp]
\centering
\caption{Sections of this chapter $\to$ in-app help anchors}
\label{tab:ch06-help}
\begin{tabular}{ll}
\hline
Section & Help topic \\
\hline
Show the lanes & \texttt{automation-show-lanes} \\
Draw and edit the curve & \texttt{automation-edit-curve} \\
Write modes & \texttt{automation-write-modes} \\
Bind a parameter & \texttt{automation-target-binding} \\
LFO boundary & \texttt{modulation-lfo-boundary} \\
\hline
\end{tabular}
\end{table}

\paragraph{What you will be able to do} Show and edit automation curves, explain Off/Write/Touch/Latch, and know that the LFO exists in the DSP core but is not yet wired into the synthesizer.

\paragraph{Hands-on exercise} Expand an automation lane on a track, select the volume parameter, use Touch mode to nudge the volume while playing and confirm only that touched segment was written; switch back to Off, drag a vertex manually and compare by ear; finally look for an LFO vibrato control and confirm it is not currently provided.
